% Copyright (c) 2026 Eugene Joseph M. Ragasa.
% SPDX-License-Identifier: Apache-2.0
% Author: Eugene Joseph M. Ragasa
% Formal proof completed: 2026-10-20
\documentclass[11pt]{article}
\usepackage[T1]{fontenc}\usepackage{amsmath,amssymb}\usepackage[margin=1in]{geometry}\usepackage[hidelinks]{hyperref}
\title{Proof Dossier: Canonical Rank-Two Band Order}\author{Eugene Joseph M. Ragasa}\date{Formal proof completed: 2026-10-20}
\begin{document}\maketitle
\section{Identity and status}
ID \texttt{BG-FOUNDATION-THM-014}; status \texttt{proved}. Lean declaration:
\par\noindent{\tiny\path{ProjectKoios.BandGauge.Foundations.}}\\
{\tiny\path{rankTwoLowerBand_le_upperBand}}.
\section{Source and contribution boundary}
Source-neutral supporting algebra; relation \texttt{background-only} to standard finite
Hermitian matrix theory through pinned Mathlib. No publication-source or novelty claim.
\section{Statement and proof}
For $r=\sqrt{x^2+y^2+z^2}$, the canonical bands satisfy $t-r\le t+r$ because
$r\ge0$.
\section{Boundary}
This orders the two exact coordinate-defined bands. It does not certify numerical
sorting or eigensolver output. Kernel axioms are \texttt{propext},
\texttt{Classical.choice}, and \texttt{Quot.sound}; no admissions. Checked with Lean \texttt{4.35.0} and Mathlib revision \texttt{56a005855ea17e0b0958f73ec801ae140b0a1679}. AI-assisted and kernel-checked. No independent expert review is recorded.
\end{document}
